\chapter{Mammary Infection}

\section*{Definition and Overview}
Mammary infection, commonly referred to as mastitis, is an inflammatory condition of the breast tissue that is frequently, though not always, accompanied by an infectious pathogen. While it is most commonly associated with lactating women (lactational mastitis), it can also occur in non-lactating females and, rarely, in males (non-lactational mastitis). The condition is characterised by localised breast pain, warmth, erythema, and swelling, often accompanied by systemic symptoms such as fever and malaise. In lactating women, it is typically caused by milk stasis and subsequent bacterial entry. If left untreated or managed inappropriately, mammary infections can progress to breast abscesses, requiring surgical intervention.

\section*{Epidemiology}
Lactational mastitis is a common complication of breastfeeding, estimated to affect approximately 10\% to 30\% of lactating mothers, with the peak incidence occurring within the first six weeks postpartum. Non-lactational mastitis is less common and is typically subdivided into periductal mastitis (commonly affecting young or middle-aged women, strongly associated with cigarette smoking) and idiopathic granulomatous mastitis (a rare chronic inflammatory condition). In Australia, breast infections are a frequent reason for postpartum general practice consultations and maternal hospital presentations.

\section*{Aetiology and Pathophysiology}
The development of mammary infections involves anatomical barriers, fluid stasis, and microbial colonisation.
\begin{itemize}
    \item \textbf{Milk stasis:} The primary predisposing factor in lactational mastitis, caused by poor latch, infrequent feeding, blocked milk ducts, or oversupply, resulting in stagnant milk that acts as a culture medium.
    \item \textbf{Bacterial entry:} Pathogens enter the breast tissue through cracked or damaged nipples or via the milk ducts. The most common causative organism is Staphylococcus aureus, followed by Streptococcus species and Escherichia coli.
    \item \textbf{Smoking and duct damage:} In non-lactational periductal mastitis, cigarette smoking damages the subareolar duct epithelium, leading to squamous metaplasia, keratin plugging, duct ectasia, and secondary anaerobic bacterial infection.
    \item \textbf{Mechanical trauma:} Direct trauma to the breast, nipple piercings, or friction from tight clothing can compromise skin integrity and facilitate bacterial access.
    \item \textbf{Immunosuppression:} Underlying conditions such as diabetes mellitus or immunosuppressive therapy increase the risk of developing both lactational and non-lactational infections.
\end{itemize}

\section*{Clinical Features}
Mammary infections present with both localised breast symptoms and systemic manifestations.
\begin{itemize}
    \item \textbf{Localised breast pain:} A sharp, burning, or aching pain, typically unilateral and localised to a single breast quadrant.
    \item \textbf{Erythema and warmth:} A wedge-shaped area of redness, swelling, and warmth that feels tender to the touch.
    \item \textbf{Systemic symptoms:} Sudden onset of high fever (often >38.5$^\circ$C), chills, rigours, body aches, and fatigue, mimicking a flu-like illness.
    \item \textbf{Nipple changes:} Cracks, fissures, or purulent discharge from the nipple.
    \item \textbf{Palpable mass:} A tender, firm area of breast tissue, which may indicate a blocked duct or, if fluctuant, an evolving breast abscess.
\end{itemize}

\section*{Differential Diagnosis}
It is critical to distinguish mammary infections from non-infectious inflammatory conditions and underlying malignancy.
\begin{itemize}
    \item \textbf{Breast abscess:} A localised collection of pus within the breast tissue, characterised by fluctuance and a lack of improvement after 48 hours of antibiotic therapy.
    \item \textbf{Inflammatory breast cancer (IBC):} A rare, highly aggressive form of breast cancer presenting with breast erythema, warmth, and skin thickening (peau d'orange) that resembles mastitis but does not resolve with antibiotics.
    \item \textbf{Blocked milk duct:} A localised, tender breast lump without systemic symptoms or significant erythema, resolving with conservative drainage measures.
    \item \textbf{Engorgement:} Bilateral, generalised, painful breast swelling occurring shortly after birth due to vascular congestion, resolving with feeding or milk expression.
    \item \textbf{Fibroadenoma or cyst:} Benign, well-circumscribed, non-tender masses that do not present with acute inflammatory or systemic signs.
\end{itemize}

\section*{When to Seek Medical Care}
Women experiencing breast swelling, redness, or pain should seek prompt medical advice from a general practitioner or lactation consultant.

\textbf{Call 000 (or local emergency services) for:}
\begin{itemize}
    \item Signs of systemic sepsis, including rapid breathing, confusion, tachycardia, dizziness, or a severe, uncontrolled high fever.
    \item Spreading redness that rapidly involves the entire breast, accompanied by severe, unmanageable pain.
\end{itemize}
Other reasons to see a doctor: If a breast lump persists for more than a few weeks after the infection clears, or if there is no improvement after 48 hours of starting antibiotics, medical evaluation is required.

\section*{Diagnosis and Investigations}
Diagnosis is primarily clinical, but investigations are indicated in severe, recurrent, or non-responding cases.
\begin{itemize}
    \item \textbf{Clinical evaluation:} Detailed physical examination of both breasts, evaluating for signs of abscess, skin changes, and systemic toxicity.
    \item \textbf{Breast milk culture:} Indicated in hospital-acquired, recurrent, or severe mastitis, or if the infection fails to respond to initial empiric antibiotics, to identify the causative organism and its sensitivity.
    \item \textbf{Breast ultrasound:} The primary imaging modality used to detect fluid collections (abscesses) and guide needle aspiration.
    \item \textbf{Mammography:} Performed in non-lactational cases or after acute inflammation resolves to rule out underlying malignancy, particularly if a breast mass persists.
    \item \textbf{Fine-needle aspiration (FNA) or biopsy:} Indicated for histological evaluation if idiopathic granulomatous mastitis is suspected or to rule out inflammatory breast cancer.
\end{itemize}

\section*{Management}
Management involves supporting breastfeeding, promoting milk drainage, and administering appropriate pharmacotherapy.
\begin{itemize}
    \item \textbf{Continued breastfeeding:} The most critical step; mothers should be encouraged to continue feeding from the affected breast to empty the breast and prevent further milk stasis.
    \item \textbf{Pharmacotherapy (Antibiotics):} Empiric oral antibiotics targeting Staphylococcus aureus (e.g.\ flucloxacillin or cephalexin) are prescribed for 5 to 10 days; if penicillin-allergic, clindamycin or erythromycin is used.
    \item \textbf{Symptomatic relief:} Non-steroidal anti-inflammatory drugs (NSAIDs, such as ibuprofen) and paracetamol to reduce pain and inflammation; warm compresses before feeding and cold compresses after feeding can help.
    \item \textbf{Lactation support:} Correcting breastfeeding technique, improving infant latch, and using breast expression (manual or pump) if the breast remains full after feeding.
    \item \textbf{Abscess drainage:} If an abscess develops, ultrasound-guided needle aspiration is preferred over surgical incision and drainage to minimise scarring and allow continued breastfeeding.
\end{itemize}

\section*{Complications}
Complications can arise from delayed treatment, poor milk drainage, or incorrect antibiotic therapy.
\begin{itemize}
    \item \textbf{Breast abscess:} The most common complication, occurring in 3\% to 11\% of mastitis cases, characterised by a localised, fluctuant collection of pus.
    \item \textbf{Recurrent mastitis:} Repeated episodes of infection, often due to unresolved milk stasis, improper latch, or inadequate antibiotic duration.
    \item \textbf{Cessation of breastfeeding:} Premature termination of breastfeeding due to severe pain, anxiety, or incorrect medical advice.
    \item \textbf{Fistula formation:} Chronic communications between the breast ducts and the skin, particularly in periductal non-lactational mastitis.
    \item \textbf{Systemic sepsis:} Rare but life-threatening spread of infection into the bloodstream, requiring intravenous antibiotics and intensive care.
\end{itemize}

\section*{Prognosis and Outcomes}
With early diagnosis and appropriate management, the prognosis for lactational mastitis is excellent. Most cases resolve completely within 2 to 5 days of starting antibiotic therapy and optimising breast drainage. Continuing to breastfeed does not harm the infant and is associated with faster recovery and lower recurrence rates. Non-lactational mammary infections, particularly periductal and granulomatous mastitis, tend to have a more chronic, relapsing course, often requiring prolonged treatment, smoking cessation, and close specialist follow-up.

\section*{Prevention and Self-Care}
Preventative measures focus on maintaining good lactation hygiene and avoiding milk stasis.
\begin{itemize}
    \item \textbf{Ensure proper positioning and latch:} Seek early assistance from a midwife or lactation consultant to ensure the infant is attaching correctly.
    \item \textbf{Feed on demand:} Breastfeed frequently and avoid skipping feedings or scheduling rigid feeding intervals.
    \item \textbf{Empty breasts fully:} Allow the infant to finish the first breast before offering the second; use a breast pump to express remaining milk if necessary.
    \item \textbf{Avoid breast compression:} Wear supportive but loose-fitting bras and avoid sleeping on the stomach to prevent blocking milk ducts.
    \item \textbf{Nipple skincare:} Keep nipples clean and dry; apply purified lanolin cream or a drop of expressed breast milk to soothe cracked nipples.
\end{itemize}

\section*{Sources and Further Reading}
The following organisations provide further information and support:
\begin{itemize}
    \item Australian Breastfeeding Association. \textit{Support services, helpline, and resources on mastitis and breastfeeding management.} https://www.breastfeeding.asn.au/
    \item Royal Women's Hospital (Melbourne). \textit{Clinical guidelines and patient brochures on mastitis and breast infections.} https://www.thewomens.org.au/
    \item Healthdirect Australia. \textit{Overview of mastitis, including symptoms, treatments, and self-care tips.} https://www.healthdirect.gov.au/mastitis
    \item NHS. \textit{Consumer health information on mastitis, symptoms, and when to seek advice.} https://www.nhs.uk/conditions/mastitis/
    \item Academy of Breastfeeding Medicine. \textit{Clinical protocols and guidelines for managing mastitis spectrum disorders.} https://www.bfmed.org/
\end{itemize}
